
\begin{frame}[fragile]{Qualité d'un schéma}

Cette session discute des bons et mauvais schémas relationnels. 

\vfill

Les notions
introduites sont:

\begin{redbullet}
  \item Les anomalies et incohérences dues à un schéma défectueux
  \item Les dépendances fonctionnelles
  \item Les clés primaires et étrangères
  \item La \alert{normalisation} 
\end{redbullet}

\vfill
\begin{blocgauche}
\alert{Ces diapositives correspondent au support en ligne disponible sur
le site \url{http://sql.bdpedia.fr/}}
\end{blocgauche}

\end{frame}


\begin{frame}{Pourquoi pas une unique table?}

\begin{small}
\begin{center}
\begin{tabular}{l|c|l|l|c}
 titre &  année &  nomMES &  prénomMES &  annéeNaiss \\
\hline
Alien & 1979 & Scott & Ridley & 1943\\
Vertigo  & 1958 & Hitchcock & Alfred & 1899\\
Psychose & 1960 & Hitchcock & Alfred & 1899\\
Kagemusha & 1980 & Kurosawa & Akira & 1910\\
Volte-face & 1997 & Woo & John & 1946\\
Pulp Fiction & 1995 & Tarantino & Quentin & 1963 \\
Titanic & 1997 & Cameron & James & 1954\\
Sacrifice & 1986 & Tarkovski & Andrei &1932\\
\hline
\end{tabular}
\end{center}
\end{small}

\vfill

Risques d'anomalies et d'incohérences en \alert{insertion}, \alert{modification}
et \alert{suppression}. 
\end{frame}


\begin{frame}{Notion de dépendance fonctionnelle}

Il y a \alert{dépendance fonctionnelle} $A \to B$ entre deux attributs $A$ et $B$ d'une relation
$R$ quand la connaissance de la valeur de $A$ implique la connaissance de la valeur de $B$.

\vfill

Exemple: le schéma \texttt{ (nom, prénom, noSS, dateNaissance, adresse, email)}

\vfill
Les dépendances:
\begin{redbullet}
  \item $email \to nom, pr\acute{e}nom, noSS, dateNaissance, adresse$
 \item $noSS \to  email, nom, pr\acute{e}nom, dateNaissance, adresse$
\end{redbullet}

\vfill

\begin{blocgauche}
Les DF servent à analyser la qualité d'un schéma.
\end{blocgauche}

\end{frame}

%%%%%%%%%%%%%%%
\begin{frame}{Les manuscrits et les experts}

Un éditeur fait évaluer tous les manuscrits qu'on lui soumet par un expert. 

\vfill 
Voici le schéma contenant tous les attributs.

\vfill

\texttt{Manuscrit (id\_manuscrit, auteur, titre, id\_expert, nom, commentaire)}

\vfill

Et les dépendances \alert{minimales} et \alert{directes}:

\vfill
\begin{redbullet}
  \item $id\_manuscrit \to auteur, titre, id\_expert, commentaire$
 \item $id\_expert \to  nom$
\end{redbullet}

\vfill

\begin{blocgauche}
Ce schéma a le même type de problème que celui des films. Nous allons l'analyser avec les DF.
\end{blocgauche}

\end{frame}


%%%%%%%%%%%%%%%
\begin{frame}{Un exemple de la relation \texttt{Manuscrit}}
 
\begin{small}
\begin{tabular}{c|l|p{4cm}|c|p{2cm}|p{3cm}}
\textbf{id\_m} & \textbf{auteur} & \textbf{titre} & \textbf{id\_exp}& \textbf{nom}& \textbf{commentaire} \\ \hline
   10& Serge&  L'arpète &2& Philippe & Une réussite\\
   20& Cécile& L'art du chant grégorien& 2& Sophie & Bravo\\
   10& Serge&  L'arpète &2& Philippe & Une réussite\\
   10& Philippe&  SQL &1& Sophie & La référence\\
\hline
\end{tabular}
\end{small}

\vfill
\begin{redbullet}
\item $id\_expert \to nom$ n'est pas respectée par le premier et deuxième
    nuplet.
    \item $id\_manuscrit \to auteur, titre, id\_expert, commentaire$ n'est pas respectée 
    par le premier et quatrième nuplet.
\end{redbullet}

\vfill

\begin{blocgauche}
Les DF sont respectées par le premier et le troisième, mais
du coup ce sont des \alert{doublons}. Il faut \alert{éviter}
ces anomalies.
\end{blocgauche}
\end{frame}

%%%%%%%%%%%%%%%
\begin{frame}{Notion de clé}


\begin{defblock}{Clé}
Une clé d'une relation $R$ est un sous-ensemble \alert{minimal} $C$
des attributs tel que tout attribut de $R$ dépend fonctionnellement de $C$. 
\end{defblock}
\vfill

\begin{redbullet}
  \item $id\_manuscrit$ est une clé pour la relation \texttt{Manuscrit}  
 \item $(id\_manuscrit, auteur)$ n'est pas une clé car non minimal
\end{redbullet}

\vfill
\begin{blocgauche}
On peut trouver plusieurs clés. On en choisit une comme \alert{clé primaire}.
Les autres sont clés secondaires.
\end{blocgauche}
\end{frame}


%%%%%%%%%%%%%%%
\begin{frame}{Normalisation}
  
\begin{defblock}{Forme normale (troisième)}
  Un schéma de relation $R$ est \alert{normalisé} quand, dans
   toute dépendance fonctionnelle $S \to A$ sur
   les attributs de $R$, $S$ est une clé. 
\end{defblock}
 
\vfill

Le schéma précédent n'est pas normalisé car il existe une dépendance fonctionelle
$id\_expert \to nom$, alors que l'attribut
$id\_expert$ n'est pas une clé. 

\vfill

\begin{blocgauche}
On peut \alert{toujours} se ramener à des relations normalisées.
\end{blocgauche}

\end{frame}


%%%%%%%%%%%%%%%
\begin{frame}{La bonne solution}
 
Voici le schéma de nos manuscrits, \alert{normalisé} 
\vfill

\begin{redbullet}
  \item \texttt{Manuscrit (id\_manuscrit,  auteur, titre, id\_expert, commentaire)}
 \item \texttt{Expert (id\_expert, nom)}
\end{redbullet}

\vfill

Elles sont toutes deux en troisième forme normale (3FN)  \alert{et} on n'a pas perdu d'information

\vfill


\begin{blocgauche}
\alert{Important}:  l'attribut \texttt{id\_expert} dans \texttt{Manuscrit} est une \alert{clé étrangère}.
 \end{blocgauche}

\end{frame}


%%%%%%%%%%%%%%%
\begin{frame}{Pour n'a-t-on pas perdu d'information?}
 
Parce qu'on peut reconstituer \alert{exactement} la relation complète
\vfill

    
\begin{redbullet}
  \item Connaissant \texttt{id\_expert}, je connais 
     toutes les autres informations sur l'expert.
 \item Connaissant \texttt{id\_manuscrit}, je connais \texttt{id\_expert}
     \item  Par \alert{transitivité}, connaissant \texttt{id\_manuscrit}, je  connais \texttt{id\_expert},
     et connaissant \texttt{id\_expert}, je connais toutes les autres informations sur l'expert
\end{redbullet}

\vfill


\begin{blocgauche}
L'association de la clé primaire et de la clé étrangère est le mécanisme
de liaison des nuplets en relationnel.
\end{blocgauche}

\end{frame}


%%%%%%%%%%%%%%%
\begin{frame}{Avec un exemple}
 
\begin{tabular}{c|c|c}
\textbf{id\_expert} & \textbf{nom} & \textbf{adresse}\\ \hline
   1 &Sophie &rue Montorgueil, Paris\\
   2& Philippe  &rue des Martyrs, Paris\\
\hline
\end{tabular}

\vfill

\begin{tabular}{c|l|p{4cm}|c|p{3cm}}
\textbf{id\_manuscrit} & \textbf{nom} & \textbf{titre} & \textbf{id\_expert}  & \textbf{Commentaire}\\ \hline
   10& Serge&  L'arpète &2& Une réussite\\
   20& Cécile& Un art du chant grégorien sous Louis XIV& 1& Bravo\\
\hline
\end{tabular}

\vfill

\begin{blocgauche}
La valeur d'une clé étrangère est \alert{toujours} celle
d'une clé de la table référencée.
\end{blocgauche}
\end{frame}


%%%%%%%%%%%%%%%
\begin{frame}{Schéma normalisé et  contraintes}
 
On normalise les schémas, et on demande au système de garantir deux contraintes.

\begin{defblock}{Contrainte d'unicité}
Une valeur de clé ne peut apparaître \alert{qu'une fois} dans une relation.
\end{defblock}
 
\vfill

\begin{defblock}{Contrainte d'intégrité référentielle}
La valeur d'une clé étrangère doit \alert{toujours}
   être également une des valeurs de la clé référencée.\end{defblock}

\vfill

\begin{blocgauche}
On obtient des schémas sans anomalie.
\end{blocgauche}

\end{frame}

%%%%%%ù

\begin{frame}{À retenir}

Un schéma de BD correct:
\vfill
\begin{redbullet}
  \item Toute relation a une clé
 \item Tous les attributs dépendent \alert{directement} de la clé (forme normale)
 \item On peut lier les nuplets les uns aux autres avec des clés étrangères.
 \item Le système évite les anomalies en garantissant
    la \alert{contrainte d'unicité} (clé) et la \alert{contrainte d'intégrité référentielle} (clé étrangère)
\end{redbullet}

\vfill

\alert{Comment trouve-t-on le bon schéma\,?} Nous y reviendrons.
\end{frame}

